\documentclass[12pt,a4paper]{article}

\usepackage[margin=2.5cm]{geometry}
\usepackage{fontspec}
\usepackage{xeCJK}
\usepackage{graphicx}
\usepackage{booktabs}
\usepackage{longtable}
\usepackage{array}
\usepackage{enumitem}
\usepackage{hyperref}
\usepackage{xcolor}
\usepackage{listings}
\usepackage{caption}
\usepackage{subcaption}
\usepackage{float}

\setmainfont{Times New Roman}
\setsansfont{Arial}
\setmonofont{Menlo}
\setCJKmainfont[
  Path=/System/Library/Fonts/Supplemental/,
  FontIndex=7,
  BoldFont=Songti.ttc,
  BoldFeatures={FontIndex=2}
]{Songti.ttc}
\setCJKsansfont[
  Path=/System/Library/Fonts/Supplemental/,
  FontIndex=7,
  BoldFont=Songti.ttc,
  BoldFeatures={FontIndex=2}
]{Songti.ttc}
\setCJKmonofont[
  Path=/System/Library/Fonts/Supplemental/,
  FontIndex=7
]{Songti.ttc}

\hypersetup{
  colorlinks=true,
  linkcolor=blue,
  urlcolor=blue,
  citecolor=blue
}

\lstdefinestyle{code}{
  basicstyle=\small\ttfamily,
  breaklines=true,
  frame=single,
  columns=fullflexible,
  tabsize=2
}
\lstset{style=code}

\newcommand{\projecttitle}{心音之忍 Heartbeat of Shinobi}
\newcommand{\coursename}{Computer Programming II}
\newcommand{\semester}{Spring 2026}

\title{
  \vspace{-1.5cm}
  \textbf{\projecttitle}\\
  \large Final Project Written Report
}
\author{
  組員一：姓名／學號 \\
  組員二：姓名／學號 \\
  組員三：姓名／學號
}
\date{2026 年 6 月 12 日}

\begin{document}
\maketitle

\begin{abstract}
本文說明本組 final project 的主題、設計動機、系統架構、核心技術、實作成果、測試方式與組內分工。本專案是一款以 Godot 與 GDExtension 實作的 2D 橫向動作遊戲 prototype，核心特色包含精準格擋、架勢值、心跳 BPM 回饋、敵人 AI、Boss 戰、道具與地圖流程。
\end{abstract}

\tableofcontents
\newpage

\section{專案概述}

\subsection{專案名稱與一句話介紹}
\begin{itemize}[leftmargin=2em]
  \item 專案名稱：\projecttitle
  \item 類型：2D 橫向動作遊戲 prototype
  \item 一句話介紹：玩家透過攻擊、格擋、閃避與道具，在心跳壓力逐漸升高的戰鬥中擊敗敵人與 Boss。
\end{itemize}

\subsection{設計動機}
請說明為什麼選擇這個題目、想解決或探索什麼玩法問題，以及本專案和一般 2D 動作遊戲相比的差異。

\subsection{最終成果摘要}
請用條列方式整理完成的功能，例如：
\begin{itemize}[leftmargin=2em]
  \item 玩家移動、跳躍、爬牆、攻擊、格擋、閃避與道具使用。
  \item 敵人巡邏、追擊、攻擊、受擊、格擋、架勢崩潰與死亡。
  \item Boss 戰、Boss 連段、階段轉換、處決流程與 HUD。
  \item 心跳 BPM 系統、畫面與音效回饋。
  \item C/C++ GDExtension 戰鬥計算與 GDScript fallback。
  \item 測試腳本、debug overlay 與 autopilot 驗證工具。
\end{itemize}

\section{遊戲內容與使用方式}

\subsection{遊戲流程}
請描述玩家從開始畫面、進入關卡、遭遇敵人、收集或使用道具、進入 Boss 戰到完成 prototype 的流程。

\subsection{操作方式}
\begin{table}[H]
\centering
\caption{主要操作}
\begin{tabular}{lll}
\toprule
功能 & 按鍵 & 說明 \\
\midrule
左右移動 & A / D 或方向鍵 & 控制角色水平移動 \\
跑步 & Shift 長按 & 提高移動速度 \\
跳躍／攀爬 & Space & 跳躍，靠牆時可攀爬 \\
攻擊 & J 或滑鼠左鍵 & 執行基本攻擊 \\
格擋／彈刀 & K 或滑鼠右鍵 & 短按嘗試 parry，長按 block \\
閃避 & L & 短距離閃避 \\
選擇道具 & 1--5 & 切換道具 \\
使用攻擊道具 & E & 使用苦無、煙霧彈等 \\
使用補血道具 & R & 使用葫蘆或補血道具 \\
\bottomrule
\end{tabular}
\end{table}

\subsection{執行環境與執行方式}
本課程規格要求 final project 必須能在 Ubuntu 24.04 環境下運作。請在此處寫明安裝需求與實際執行步驟。

\begin{lstlisting}[language=bash,caption={範例建置與執行指令，請依實際情況修正}]
sudo apt update
sudo apt install -y build-essential scons pkg-config \
  libx11-dev libxcursor-dev libxinerama-dev libxrandr-dev \
  libxi-dev libgl-dev libasound2-dev libpulse-dev

scons platform=linux target=template_debug arch=x86_64 api_version=4.6
godot --editor project/
\end{lstlisting}

\section{系統架構與技術實作}

\subsection{整體架構}
請放入架構圖或以文字說明各層關係。建議包含：
\begin{itemize}[leftmargin=2em]
  \item Godot scene layer：場景、角色、地圖、UI 與音效。
  \item GDScript gameplay layer：玩家、敵人、Boss、道具與流程控制。
  \item Combat facade layer：\texttt{combat\_math.gd}、\texttt{combat\_server.gd}、\texttt{combat\_runtime.gd}。
  \item Native C/C++ layer：\texttt{combat\_math.c}、\texttt{combat\_mgr.c} 與 GDExtension wrapper。
\end{itemize}

\subsection{核心玩法一：精準格擋與架勢值}
請說明 parry / block 的判定方式、輸入緩衝、攻擊 active frame、架勢值累積與處決條件。

\subsection{核心玩法二：心跳 BPM 系統}
請說明心跳如何隨戰鬥狀態改變，如何影響格擋容錯、架勢恢復、傷害或視覺音效回饋。

\subsection{核心玩法三：敵人與 Boss AI}
請說明一般敵人與 Boss 的狀態機、攻擊選擇、追擊與保持距離、連段、階段轉換或反制邏輯。

\subsection{C/C++ 與 GDExtension 使用}
本課程規格要求整個 project 至少有一部分使用 C 語言。請在此處說明：
\begin{itemize}[leftmargin=2em]
  \item 哪些模組使用 C 或 C++ 實作。
  \item 為什麼這些邏輯適合放在 native layer。
  \item Godot 如何透過 GDExtension 呼叫 native 函式。
  \item 若 native library 未成功載入，GDScript fallback 如何維持可執行性。
\end{itemize}

\section{創意、藝術與設計}

\subsection{創意與原創性}
請說明本 project 的創新點。可從玩法、技術、敘事、視覺、回饋設計或開發工具等角度描述。

\subsection{美術、音效與 UI}
請描述使用的 sprite、地圖素材、音效、BGM、HUD、debug overlay、畫面震動、心跳回饋等設計。

\subsection{關卡與體驗設計}
請說明關卡如何安排敵人、道具、checkpoint、Boss 戰與難度曲線。

\section{測試與可執行性}

\subsection{測試方法}
請說明使用哪些方式確認功能正確，例如人工測試、Godot 測試腳本、C 單元測試、debug overlay、autopilot。

\begin{table}[H]
\centering
\caption{測試項目摘要}
\begin{tabular}{p{0.28\linewidth}p{0.42\linewidth}p{0.2\linewidth}}
\toprule
測試項目 & 測試內容 & 結果 \\
\midrule
玩家移動與攻擊 & 測試走路、跑步、跳躍、攻擊連段 & 待填 \\
格擋與 parry & 測試 parry window、block、架勢值變化 & 待填 \\
Boss 戰 & 測試 Boss 攻擊、階段轉換、處決流程 & 待填 \\
GDExtension & 測試 native library 載入與 fallback & 待填 \\
Ubuntu 24.04 執行 & 測試建置、啟動與操作流程 & 待填 \\
\bottomrule
\end{tabular}
\end{table}

\subsection{已知限制}
請列出目前仍未完成或受時間限制的部分，例如素材完整度、平衡性、關卡長度、平台相容性、效能或測試覆蓋率。

\section{組內分工}

請依實際貢獻填寫。課程規格要求書面報告必須包含組內分工；若分工差異過大，分數可能依個人狀況調整。

\begin{longtable}{p{0.17\linewidth}p{0.18\linewidth}p{0.42\linewidth}p{0.13\linewidth}}
\caption{組員分工表}\\
\toprule
組員 & 學號 & 主要負責內容 & 貢獻比例 \\
\midrule
\endfirsthead
\toprule
組員 & 學號 & 主要負責內容 & 貢獻比例 \\
\midrule
\endhead
姓名一 & 學號一 & 待填：例如玩家控制、戰鬥系統、C/GDExtension、測試等。 & 待填 \\
姓名二 & 學號二 & 待填：例如敵人 AI、Boss 戰、關卡、UI、音效等。 & 待填 \\
姓名三 & 學號三 & 待填：例如美術素材整合、道具、報告、README、測試等。 & 待填 \\
\bottomrule
\end{longtable}

\section{LLM 使用揭露}

課程規格要求若使用 LLM，必須在報告中明確告知使用範圍。請依實際情況保留或刪改以下內容。

\begin{table}[H]
\centering
\caption{LLM 使用紀錄}
\begin{tabular}{p{0.24\linewidth}p{0.34\linewidth}p{0.32\linewidth}}
\toprule
使用範圍 & 具體用途 & 人工檢查與修改 \\
\midrule
程式碼 & 待填：例如輔助撰寫測試、整理 fallback 邏輯、debug 建議。 & 待填：人工閱讀、測試、修改後才納入。 \\
報告文字 & 待填：例如協助整理章節架構或潤稿。 & 待填：由組員確認內容符合實際實作。 \\
README / 文件 & 待填：例如協助整理執行步驟。 & 待填：以實際可執行結果校正。 \\
\bottomrule
\end{tabular}
\end{table}

\subsection{本專案使用 LLM 的方式}

本專案開發過程中使用 ChatGPT / Codex 與 Gemini 作為輔助工具。主要用途包含：理解既有程式碼、規劃功能實作、協助撰寫或修改 Godot / GDScript / C / C++ 程式碼、排查錯誤、整理 README 與報告模板、討論遊戲系統設計，以及檢查專案狀態。LLM 產生的內容均由組員人工檢查，並以實際可執行結果、遊戲內測試與版本控制紀錄確認後才納入專案。

\subsection{楊立宇使用 Codex CLI 的 Prompt 紀錄}

\begin{lstlisting}[caption={楊立宇使用 Codex CLI 的 prompt 紀錄}]
## 2026-05-27


1. 可以幫我裝上 https://github.com/Coding-Solo/godot-mcp.git 這個mcp嗎

2. 我該怎麼用godot mcp

3. 可以請你幫我確認現在素材的狀態嗎

## 2026-05-31

1. 請幫我到feature/MapCollision上，做出角色爬牆的功能，如果缺少素材的話請直接告訴我

2. 請先幫我push上去，並在commit message 中標註這是mvp等級的爬牆

3. 請幫我加上邊界判斷，如果掉出世界就判斷為死亡

4. 我想在這遊戲加入存檔的功能，你覺得可行嗎

5. 應該說我想的是像魂系遊戲一樣，有個火點，坐下去怪物就會全部重生，所以不需要在意敵人血量，只需要紀錄玩家的狀態就好

6. 現在有checkpoint的素材嗎，請幫我檢查一下

## 2026-06-03

1. 目前想做的道具系統應該是長這樣的
   血瓶擴容
   腎上腺素 -> 升心跳
   血壓藥 -> 降心跳
   煙霧彈 -> 暫停boss動作
   苦無 -> 飛雷神（丟出後，按下某個按鈕，瞬間移動到苦無當下位置）

2. 可以幫我把丟出去的苦無變成 kunai_throwed.png嗎

3. 請幫我將圖片轉180度，再縮小1/3

4. 請再小1/2

5. ash_ball也是一樣，丟出去的幫我換成 @project/assets/ash_balls/ash_ball_throwed.png

6. 請幫我調整角色的攻擊速度，讓他快大概0.3秒左右，但攻擊完需要大概0.3秒的硬直

7. 先幫我 push上去

8. 可以幫我把ash ball的 smoke特效對齊一點嗎。ash ball丟出去後碰到地板為終止嗎

9. 請幫我push上去

10. 請幫我看一下現在的狀態該不該留著，如果boss的血量有跟軀幹值有點問題的話，幫我開新分支處理這個問題

11. 可以幫我把boss的血量跟軀幹值的數字移除嗎，另外按L會一直往前，即便一邊按下A也會一直往前，可以幫我修一下嗎

12. 那先把feature/item-system合併進main好了

## 2026-06-05

1. 請幫我開新分支，優化boss的戰鬥邏輯

2. 目前最新的branch 似乎在 origin/demo-polish上 我怕直接從main開始做會有衝突，可以幫我看看嗎

3. 好的，那就麻煩你優化boss的戰鬥邏輯了

4. 可以幫我修正一個問題嗎，目前掉出地圖死亡後（顯示「心音斷絕」的那種），重新開始後會讀到掉出地圖前的血量狀態，但我覺得應該血量是滿的才正常

5. 請幫我push 上去，並幫我考慮要不要合併到demo-polish上

6. 那請幫我merge


## 2026-06-08
1. 有什麼可能會觸發bug的edge case嗎

2. savegame.json在哪裡

## 2026-06-10

1. 請幫我閱讀 "final project spec.pdf" ，並做出書面報告的模板 (.tex)，分工表先預留三個人的位置。
\end{lstlisting}

\subsection{楊立宇使用 gemini-cli 的 Prompt 紀錄}

\begin{lstlisting}[caption={楊立宇使用 gemini-cli 的 prompt 紀錄}]
1. 專案啟動與基礎建設 (2026-05-11)
  這階段主要集中在理解現有程式碼、重構建置腳本以及撰寫文件。
   * please help me understand the code base
   * help me check if SConstruct is available for All platform(especially Ubuntu20.04)
   * please help me write README.md

2. 戰鬥系統與「危」字攻擊開發 (2026-05-16)
  開始實作核心戰鬥邏輯，特別是參考《隻狼》的「危」字攻擊機制。
   * 如果你要實作「危」字攻擊，你會怎麼做
   * 同意，請你幫我實作
   * 我在perry 的時候報錯了 ... Nonexistent function 'is_parry_successful'
   * 但現在沒有顯示「危」字，實際上我也看不出跟改之前有什麼差別，可以跟我說是哪一段程式負責的嗎

3. Boss 軀幹值 (Posture) 與 AI 行為 (2026-05-16 ~ 05-17)
  實作 Boss 的軀幹值系統、格擋頻率與攻擊連段。
   * 你可以幫我做boss 的軀幹值嗎(boss posture)
   * 但目前就算彈反成功，在Boss Posture那邊也沒有任何動靜
   * 可以請你幫我提高boss格擋的頻率嗎，在我攻擊的時候
   * 那你能幫我增加boss的攻擊頻率，並安排攻擊連段中間的間隔嗎

4. Boss 房實作與 UI/音效優化 (2026-05-17)
  建立獨立的 Boss 戰場景，並處理死亡音效與 UI 顯示錯誤。
   * 那可以幫我做一個boss房嗎 ... 只要一進到boss房，boss就會開始追著玩家
   * 現在是不是玩家死了之後，遊戲的背景音樂會和死亡的音樂重疊
   * 可以幫我做成死亡時，背景音樂淡出，死亡音樂淡入嗎

5. 美術素材規劃與道具系統 (2026-05-18)
  列出遊戲所需的所有美術資源，並規劃後續的道具系統。
   * 可以幫我列出所有需要的美術或圖片嗎，這個遊戲所有需要的，若已經有了就在旁邊註記，可以讀 @spec.md
   * 另外我們還有規劃要加上道具：血瓶擴容、腎上腺素、血壓藥 (降心跳)、煙霧彈 (暫停boss動作)、苦無 (飛雷神)
\end{lstlisting}

\section{結論}

請總結本 project 完成了什麼、最重要的技術或設計挑戰、學到的內容，以及如果有更多時間會如何延伸。

\appendix
\section{附錄：專案檔案結構}

\begin{lstlisting}[caption={可依實際目錄修正}]
src/                         C/C++ GDExtension source
project/                     Godot project
project/scenes/              Scenes and prefabs
project/scripts/             GDScript gameplay logic
project/assets/              Sprites, audio, maps, UI assets
project/tests/               Godot test scripts
tests/                       C unit tests
README.md                    Build and usage guide
\end{lstlisting}

\section{附錄：畫面截圖}
請放入遊戲主畫面、戰鬥畫面、Boss 戰、debug overlay 或其他重要畫面。

% \begin{figure}[H]
% \centering
% \includegraphics[width=0.85\linewidth]{path/to/screenshot.png}
% \caption{遊戲畫面截圖}
% \end{figure}

\end{document}
